\documentclass[a4paper,10pt]{report}\input{../0.Préambule/Préambule_cours_prof}\begin{document}

\definecolor{titlecolor}{rgb}{0.98, 0.4, 0.0}
\newpage
\setcounter{chapter}{7}
\chapter{Puissance de dix}

\vspace{-3mm}
\section{Définition}

\vspace{-3mm}
\begin{dft}[Puissance à exposant positif]
Si $n$ un entier naturel, alors $10^n$, qui se lit \og $10$ puissance $n$ \fg{} ou \og $10$ exposant $n$ \fg{}, est le nombre \[ 10^n=\underset{n\ \text{facteurs égaux à } 10 }{\underbrace{10\times 10 \times \dots \times 10}}=\underset{n\ \text{zéros} }{\underbrace{100 \dots 00}} \]

\vspace{-2mm}
\noindent Par convention : $10^1 = 10$ \quad et \quad $10^0 = 1$;
\end{dft}

\begin{dft}[Puissance à exposant négatif]
Soit $n$ un nombre strictement positif. On appelle « $10$ puissance moins $n$ » le nombre noté $10^{-n}$ tel que :  \[ 10^{-n}=\underset{n\ \text{zéros} }{\underbrace{0,00 \dots 01}} \]
\end{dft}

\vspace{-3mm}
\section{Les préfixes}

\vspace{-2mm}
\begin{center}
\renewcommand{\arraystretch}{1.5}
\begin{tabularx}{\linewidth}{|>{\columncolor{lightgray!80} }p{2.5cm}|*{7}{>{\centering \arraybackslash }X|}}
\hline
Préfixes & giga & méga & kilo & ... & milli & micro & nano \\\hline
Symbole & $G$ & $M$ & $k$ & ... & $m$ & $\mu$ & $n$ \\\hline
Signification & $10^{9}$ & $10^{6}$ & $10^{3}$ & ... & $10^{-3}$ & $10^{-6}$ & $10^{-9}$ \\\hline
\end{tabularx}
\end{center}

\vspace{-4mm}
\begin{ex}
1 gigaoctet, noté Go, est égal à $10^{9}$ octets, \quad ou encore $1 \mu m = 10^{-6} m$
\end{ex}

\begin{prop}
Pour multiplier un nombre décimal par $10^{n}$, on décale la virgule de $n$ rangs vers la droite.

\medskip
\noindent Pour multiplier un nombre décimal par $10^{-n}$, on décale la virgule de $n$ rangs vers la gauche.
\end{prop}

\begin{meth}

\vspace{-2mm}
\noindent \begin{minipage}{4cm}
$$A =  9,4\times 10^{3}$$

On déplace la virgule dans $9,4$ de $3$ rangs vers la droite.

$$A  = 9 400$$
\end{minipage}
\hspace{1mm}
\vline
\hspace{1mm}
\begin{minipage}{4cm}
$$B = 54\times 10^{8}$$

On déplace la virgule dans $54$ de $8$ rangs vers la droite.

$$B = 5 400 000 000$$
\end{minipage}
\hspace{1mm}
\vline
\hspace{1mm}
\begin{minipage}{4cm}
$$C = 185,003\times 10^{-2}$$

On déplace la virgule dans $185,003$ de $2$ rangs vers la gauche.

$$C = 1,85003$$
\end{minipage}
\hspace{1mm}
\vline
\hspace{1mm}
\begin{minipage}{4cm}
$$B = 54\times 10^{-8}$$

On déplace la virgule dans $54$ de $8$ rangs vers la gauche.

$$B = 0,00000054$$
\end{minipage}
\end{meth}

\vspace{-4mm}
\section{L'écriture scientifique d'un nombre décimal}

\vspace{-3mm}
\begin{dft}
Un nombre positif est écrit en notation scientifique quand il est écrit sous la forme $a \times 10^n$ avec : 
\begin{enumerate}
\item[•]$a$ un nombre décimal ne comportant qu'un seul chiffre non nul devant la virgule
\item[•]$n$ un nombre entier relatif
\end{enumerate} 

\end{dft}

\begin{ex}
L'écriture scientifique de $2600$ est $2,6 \times 10^{3}$. De même, $2600 = 2,6 \times 10^3$ ou encore $0,0137 = 1,37 \times 10^{-2}$
\end{ex}

\end{document}